\chapter{สรุปผล อภิปรายผล และข้อเสนอแนะ}

\section{สรุปผลการวิจัย}
การวิจัยนี้มีวัตถุประสงค์เพื่อ (1) ประเมินผลกระทบของโครงการ Smart Farmer ต่อการพัฒนาทักษะของเกษตรกร (2) วิเคราะห์ความสัมพันธ์ระหว่างปัจจัยที่เกี่ยวข้องกับการเปลี่ยนแปลงทักษะ และ (3) เสนอกรอบแนวทางการพัฒนาทักษะเกษตรกรโดยใช้เทคโนโลยี Machine Learning

ทั้งนี้ การสรุปผลในบทนี้จะจัดเรียงให้ตอบวัตถุประสงค์ดังกล่าวโดยเชื่อมโยงกับผลการวิเคราะห์เชิงประจักษ์ในบทที่ 4 อย่างเป็นระบบ

เพื่อให้การสรุปผลสอดคล้องกับกรอบการวิจัย บทนี้จึงเรียบเรียงข้อค้นพบเป็น 3 ส่วนตามวัตถุประสงค์การวิจัย และใช้ผลการวิเคราะห์ในบทที่ 4 เป็นฐานในการอธิบายและเชื่อมโยงประเด็นอย่างต่อเนื่อง ดังนี้

\begin{enumerate}

\item \textbf{ผลกระทบของโครงการ Smart Farmer ต่อการพัฒนาทักษะของเกษตรกร}

จากการประเมินผลกระทบด้วยวิธี Propensity Score Matching ร่วมกับ Difference-in-Differences (PSM-DID) ในบทที่ 4 พบว่า โดยภาพรวมการเข้าร่วมโครงการไม่ทำให้คะแนนทักษะของเกษตรกรเพิ่มขึ้นอย่างมีนัยสำคัญทางสถิติเมื่อพิจารณาในระดับค่าเฉลี่ย

ผลดังกล่าวแสดงให้เห็นว่า ผลลัพธ์ของโครงการอาจไม่เด่นชัดในภาพรวมภายใต้บริบทที่ผู้เข้าร่วมมีความแตกต่างกันสูง รวมถึงมีข้อจำกัดด้านการนำความรู้ไปประยุกต์ใช้ในภาคปฏิบัติ

อย่างไรก็ดี ผลการประเมินเชิงสาเหตุด้วย PSM-DID ยังมีความสำคัญในฐานะหลักฐานที่ชี้ว่า เมื่อควบคุมความแตกต่างเชิงลักษณะพื้นฐานระหว่างกลุ่มผู้เข้าร่วมและกลุ่มเปรียบเทียบแล้ว ผลสัมฤทธิ์ด้านทักษะโดยเฉลี่ยยังไม่ปรากฏเด่นชัด

\item \textbf{ความสัมพันธ์ระหว่างปัจจัยที่เกี่ยวข้องกับการเปลี่ยนแปลงทักษะ}

การวิเคราะห์เพิ่มเติมด้วยเทคนิค Machine Learning และการอธิบายตัวแบบด้วย SHAP พบว่า ปัจจัยที่สัมพันธ์กับการเปลี่ยนแปลงทักษะมีลักษณะเป็น “ปัจจัยร่วม” ที่เชื่อมโยงทั้งความพร้อมของเกษตรกรและเงื่อนไขแวดล้อมของพื้นที่

โดยปัจจัยสำคัญประกอบด้วย

\begin{itemize}
\item ระดับทักษะเดิม (baseline skills)
\item ทักษะด้านเทคโนโลยีและระดับการศึกษา
\item โครงสร้างพื้นฐานที่เอื้อต่อการผลิต เช่น ระบบชลประทาน
\end{itemize}

ข้อค้นพบนี้ช่วยขยายความหมายของผล PSM-DID โดยชี้ว่า แม้ผลกระทบเฉลี่ยไม่เด่นชัด แต่ผลลัพธ์มีความแตกต่างกันตามลักษณะของเกษตรกรและบริบทพื้นที่

\item \textbf{กรอบแนวทางการพัฒนาทักษะเกษตรกรโดยใช้ Machine Learning}

เมื่อพิจารณาผลการอธิบายตัวแบบ (SHAP) ร่วมกับผลการจัดกลุ่ม (clustering) พบว่า เกษตรกรสามารถจำแนกออกเป็นหลายกลุ่มที่มีศักยภาพและข้อจำกัดแตกต่างกันอย่างชัดเจน

จึงนำไปสู่ข้อเสนอเชิงกรอบแนวทางว่า การพัฒนาทักษะไม่ควรใช้แนวทางแบบเดียวกับทุกกลุ่ม (one-size-fits-all) แต่ควรออกแบบการสนับสนุนแบบจำเพาะกลุ่ม (targeted intervention)

โดยใช้ข้อมูลเป็นฐานเพื่อ

\begin{itemize}
\item คัดกรอง/จัดกลุ่มเกษตรกรตามความพร้อมและข้อจำกัด
\item จัดสรรรูปแบบการอบรมและการติดตามผลให้เหมาะสมกับแต่ละกลุ่ม
\item ใช้โมเดล Machine Learning เพื่อคาดการณ์โอกาสเกิดผลลัพธ์และสนับสนุนการตัดสินใจเชิงนโยบาย
\end{itemize}

\end{enumerate}

\section{อภิปรายผลการวิจัย (Discussion)}

\subsection{การตีความผลการประเมินผลกระทบ (DID)}

หัวข้อนี้มุ่งอธิบายความหมายของข้อค้นพบในบทที่ 4 โดยเชื่อมโยงผลการประเมินเชิงสาเหตุด้วย PSM-DID เข้ากับผลการวิเคราะห์เชิงลึกด้วย Machine Learning/SHAP และผลการจัดกลุ่ม (clustering) เพื่อทำความเข้าใจว่าเหตุใดผลกระทบเฉลี่ยจึงไม่เด่นชัด และเงื่อนไขใดทำให้บางกลุ่มได้รับประโยชน์มากกว่ากลุ่มอื่น
จากผลดังกล่าว สามารถอภิปรายประเด็นหลักได้ 2 ประการ ได้แก่ (1) ความเอนเอียงจากการคัดเลือกเข้าร่วมโครงการ (selection bias) และ (2) ข้อจำกัดของการสรุปผลในระดับค่าเฉลี่ยภายใต้ผลลัพธ์ที่ไม่เป็นเนื้อเดียวกัน (heterogeneous effects)

\textbf{(1) ปัญหา Selection Bias}

ผลการเปรียบเทียบลักษณะพื้นฐานในบทที่ 4 บ่งชี้ว่า เกษตรกรที่เข้าร่วมโครงการมีแนวโน้มเป็นกลุ่มที่มีความพร้อมสูงกว่าอยู่แล้ว เช่น มีทักษะเดิมสูงกว่า หรือมีเงื่อนไขแวดล้อมเอื้ออำนวยมากกว่า
แม้วิธี PSM จะช่วยจับคู่ให้กลุ่มเปรียบเทียบมีความใกล้เคียงกันมากขึ้น แต่ในทางปฏิบัติความแตกต่างที่ไม่สามารถสังเกตได้อาจยังคงอยู่

\textbf{(2) ข้อจำกัดของการวัดผลในระดับค่าเฉลี่ย}

เมื่อผลลัพธ์มีความแตกต่างกันระหว่างกลุ่ม การใช้ค่าเฉลี่ยเพียงค่าเดียวอาจทำให้ผลกระทบที่เกิดกับบางกลุ่มถูกกลบ
ดังนั้น การใช้ Machine Learning และ SHAP ในบทที่ 4 จึงช่วยเปิดเผยความแตกต่างของผลลัพธ์ในระดับรายบุคคลได้ชัดเจนยิ่งขึ้น

\subsection{บทบาทของปัจจัยโครงสร้างและทักษะเดิม}

ผลการวิเคราะห์ด้วย Machine Learning และ SHAP ชี้ว่า ปัจจัยพื้นฐานของเกษตรกรและเงื่อนไขเชิงโครงสร้างมีอิทธิพลต่อผลลัพธ์อย่างชัดเจน
โดยเฉพาะระดับทักษะเดิม ระดับการศึกษา และการเข้าถึงทรัพยากร
ข้อค้นพบนี้สะท้อนว่า การอบรมเพียงอย่างเดียวอาจไม่เพียงพอ หากไม่เสริมด้วยปัจจัยเอื้อ

\subsection{นัยเชิงระบบ (Systemic Insight)}
เมื่อพิจารณาผลทั้งหมดร่วมกันสามารถสรุปเชิงระบบได้ว่า ข้อจำกัดของผลลัพธ์โครงการไม่ได้อยู่ที่ “เนื้อหาการอบรม” เพียงอย่างเดียว แต่เกี่ยวข้องกับการออกแบบมาตรการและกลไกสนับสนุนที่ยังไม่สอดรับกับความแตกต่างของผู้เข้าร่วม ทั้งในมิติทักษะเดิม ความสามารถด้านเทคโนโลยี และข้อจำกัดเชิงโครงสร้างของพื้นที่ ดังนั้น การพัฒนาโครงการในระยะต่อไปควรมุ่งสร้างรูปแบบการสนับสนุนที่เหมาะกับแต่ละกลุ่มเป้าหมายและเพิ่มเงื่อนไขเอื้อเพื่อให้เกิดการนำความรู้ไปใช้จริง
จากการอภิปรายข้างต้น ชี้ให้เห็นว่าการยกระดับประสิทธิผลของโครงการจำเป็นต้องปรับแนวทางจากการให้บริการแบบเหมารวม ไปสู่การจัดสรรการสนับสนุนแบบจำเพาะกลุ่มบนฐานข้อมูล ซึ่งจะนำไปสู่ข้อเสนอแนะเชิงนโยบายในหัวข้อถัดไป

\section{ข้อเสนอแนะเชิงนโยบาย (Policy Recommendations)}
จากผลการศึกษา สามารถเสนอแนวทางการพัฒนาในระดับนโยบายได้ดังนี้

\subsection{การคัดกรองและจัดกลุ่มเป้าหมายบนฐานข้อมูล}
จากข้อค้นพบเรื่องความไม่เป็นเนื้อเดียวกันของผลลัพธ์ (heterogeneous effects) และบทบาทของปัจจัยพื้นฐานที่ตัวแบบ SHAP ระบุไว้ หน่วยงานผู้ดำเนินโครงการควรกำหนดระบบคัดกรองและจัดกลุ่มเกษตรกรก่อนเข้าร่วม โดยใช้ข้อมูลพื้นฐาน (เช่น ทักษะเดิม การศึกษา ทักษะเทคโนโลยี การเข้าถึงโครงสร้างพื้นฐาน และบริบทพื้นที่) เพื่อจำแนกกลุ่มความพร้อม/ข้อจำกัดอย่างน้อยในระดับ “พร้อมสูง–พร้อมปานกลาง–ข้อจำกัดสูง” การจัดกลุ่มดังกล่าวช่วยลดความเสี่ยงของการใช้ทรัพยากรกับกลุ่มที่ยังไม่พร้อมรับการอบรมขั้นสูง และช่วยเพิ่มโอกาสให้กลุ่มที่มีศักยภาพได้รับการสนับสนุนที่เหมาะสม
ในเชิงเครื่องมือ สามารถพัฒนาแบบจำลอง Machine Learning เพื่อคาดการณ์ “ความน่าจะเป็นที่จะเกิดการพัฒนาทักษะ” หลังเข้าร่วมโครงการ โดยใช้ข้อมูลย้อนหลังและตัวแปรตามที่บทที่ 4 แสดงว่าเกี่ยวข้องกับผลลัพธ์ จากนั้นนำผลการคาดการณ์ไปใช้ประกอบการคัดเลือก การจับคู่บริการสนับสนุน และการติดตามผลแบบเข้มข้นในกลุ่มที่มีความเสี่ยงไม่สามารถนำความรู้ไปใช้จริง ทั้งนี้ ควรกำหนดหลักเกณฑ์ความโปร่งใสและการคุ้มครองข้อมูลส่วนบุคคลควบคู่กันไป

\subsection{การออกแบบหลักสูตรแบบจำเพาะกลุ่ม}

\begin{itemize}
\item กลุ่มพร้อมสูง: จัดหลักสูตรขั้นสูงที่เน้นการใช้ข้อมูล/เทคโนโลยีในการตัดสินใจ (เช่น การใช้แอป/แพลตฟอร์ม การวิเคราะห์ข้อมูลการผลิต) พร้อมกิจกรรมภาคปฏิบัติและที่ปรึกษาเฉพาะด้าน เพื่อเร่งการนำไปใช้จริง
\item กลุ่มพร้อมปานกลาง: เน้นการยกระดับทักษะดิจิทัลพื้นฐาน การเข้าถึงข้อมูล และการแก้ปัญหาหน้างาน โดยเพิ่มการโค้ช/ติดตามผลหลังอบรม เพื่อช่วยข้าม “ช่องว่างการประยุกต์ใช้”
\item กลุ่มข้อจำกัดสูง: ควรเริ่มจากการเสริมทักษะพื้นฐานและการเข้าถึงปัจจัยเอื้อ (เช่น อุปกรณ์ ช่องทางสื่อสาร หรือการเชื่อมต่อบริการภาครัฐ) ก่อนเข้าหลักสูตรที่ซับซ้อน เพื่อให้การอบรมไม่กลายเป็นภาระและเพิ่มโอกาสเกิดผลลัพธ์
\end{itemize}

\subsection{การพัฒนาโครงสร้างพื้นฐาน}

ผลการวิเคราะห์ในบทที่ 4 ชี้ว่าโครงสร้างพื้นฐานและการเข้าถึงทรัพยากรมีบทบาทต่อผลลัพธ์อย่างมีนัยสำคัญ ดังนั้น ภาครัฐควรดำเนินมาตรการเสริมควบคู่การอบรม ได้แก่ (ก) การพัฒนาระบบชลประทานและการจัดการน้ำในพื้นที่เป้าหมาย (ข) การเพิ่มการเข้าถึงเทคโนโลยีและเครือข่ายสื่อสารที่จำเป็นต่อการใช้เครื่องมือดิจิทัล และ (ค) การเชื่อมโยงบริการสนับสนุนอื่น (เช่น แหล่งทุน องค์ความรู้ หรือที่ปรึกษา) เพื่อให้เกิดการนำความรู้ไปใช้จริงอย่างต่อเนื่อง

\subsection{การติดตามผลระยะยาว}
เพื่อให้การประเมินผลสะท้อนการเปลี่ยนแปลงที่เกิดขึ้นจริง ควรกำหนดระบบติดตามผลหลังการอบรมในระยะกลางและระยะยาว โดยเก็บข้อมูลก่อน–หลังเข้าร่วมอย่างเป็นมาตรฐาน (baseline และ follow-up) ทั้งในมิติทักษะ (เช่น การใช้เทคโนโลยี การจัดการข้อมูล) และมิติการนำไปใช้/ผลลัพธ์ทางการผลิต (เช่น การปรับเปลี่ยนวิธีการผลิต ต้นทุน ผลผลิต หรือการเข้าถึงตลาด) การมีข้อมูลติดตามดังกล่าวจะช่วยให้สามารถประเมินผลด้วยวิธีเชิงสาเหตุ เช่น DID ได้แม่นยำขึ้น และช่วยออกแบบการสนับสนุนเพิ่มเติมให้เหมาะกับกลุ่มที่ผลลัพธ์ยังไม่เกิด

\section{ข้อเสนอแนะสำหรับการวิจัยในอนาคต (Future Research)}

\subsection{การศึกษาระยะยาวเพื่อยืนยันผลกระทบ (Longitudinal Study)}

ควรมีการติดตามเกษตรกรในระยะยาวมากกว่าหนึ่งรอบการผลิต เพื่อประเมินว่าการอบรมส่งผลต่อการเปลี่ยนพฤติกรรมและการยกระดับทักษะอย่างยั่งยืนหรือไม่ เนื่องจากผลของการเรียนรู้และการปรับตัวมักเกิดแบบค่อยเป็นค่อยไป และอาจไม่สะท้อนชัดในช่วงเวลาสั้น ซึ่งสอดคล้องกับข้อค้นพบในบทที่ 4 ที่ผลกระทบเฉลี่ยยังไม่เด่นชัด
\subsection{การออกแบบการประเมินแบบกึ่งทดลอง/เชิงทดลองเพื่อลด Selection Bias}
งานวิจัยในอนาคตควรพิจารณาการออกแบบการประเมินที่ช่วยลดความเอนเอียงจากการคัดเลือกเข้าร่วม เช่น การสุ่มกลุ่มเข้าร่วม (Randomized Controlled Trial: RCT) ในบางพื้นที่นำร่อง หรือการใช้การออกแบบกึ่งทดลองที่เข้มแข็ง (เช่น การทยอยเปิดโครงการแบบเป็นช่วงเวลา หรือการใช้เกณฑ์คะแนน/คุณสมบัติเป็นจุดตัด) เพื่อให้สามารถระบุผลกระทบเชิงสาเหตุได้ชัดเจนขึ้นและเปรียบเทียบประสิทธิผลของรูปแบบการสนับสนุนที่แตกต่างกันไ

\subsection{การพัฒนาระบบแนะนำและการนำตัวแบบไปใช้จริง (Recommender System and Deployment)}
ควรพัฒนาระบบแนะนำ (recommender system) เพื่อจับคู่เกษตรกรกับหลักสูตร/รูปแบบการสนับสนุนที่เหมาะสมรายบุคคล โดยอาศัยข้อมูลพื้นฐานและประวัติการเข้าร่วมกิจกรรม รวมถึงการออกแบบแดชบอร์ดสำหรับเจ้าหน้าที่ภาคสนามเพื่อใช้ติดตามความก้าวหน้า ทั้งนี้ การนำตัวแบบไปใช้จริงควรพิจารณาความโปร่งใสในการตัดสินใจ ความเท่าเทียมระหว่างกลุ่มเป้าหมาย และกระบวนการทบทวน/ปรับปรุงตัวแบบอย่างสม่ำเสมอ เพื่อป้องกันการตัดสิทธิ์หรือการจัดสรรทรัพยากรที่ไม่เป็นธรรม

\subsection{การบูรณาการข้อมูลเชิงพื้นที่และบริบท (Spatial and Contextual Data)}
ควรบูรณาการข้อมูลเชิงพื้นที่และบริบท เช่น สภาพภูมิอากาศ ลักษณะพื้นที่เพาะปลูก ประเภทพืช ความเสี่ยงภัยพิบัติ รวมถึงความพร้อมของโครงสร้างพื้นฐานในพื้นที่ เพื่อเพิ่มความแม่นยำในการอธิบาย/คาดการณ์ผลลัพธ์และทำความเข้าใจความแตกต่างระหว่างกลุ่มเกษตรกร เนื่องจากปัจจัยเหล่านี้อาจเป็นตัวกำหนดความสามารถในการนำความรู้ไปใช้จริงและเป็นแหล่งที่มาของความไม่เป็นเนื้อเดียวกันของผลลัพธ์
\section{ข้อสรุปเชิงนโยบาย (Final Insight)}
จากผลการวิจัยทั้งหมดสามารถสรุปเป็นข้อสรุปเชิงนโยบายได้ว่า การพัฒนาเกษตรกรในยุคดิจิทัลควรปรับจากการฝึกอบรมแบบทั่วไป ไปสู่การพัฒนาแบบจำเพาะกลุ่มและรายบุคคลบนฐานข้อมูล โดยใช้ผลการประเมินเชิงสาเหตุ (PSM-DID) เป็นหลักฐานภาพรวม และใช้การวิเคราะห์เชิงลึกด้วย Machine Learning/SHAP ร่วมกับการจัดกลุ่มเพื่อระบุ “ใครควรได้รับการสนับสนุนแบบใด” และ “เงื่อนไขใดต้องเสริม” เพื่อให้เกิดผลลัพธ์ที่วัดได้จริง
ในทางปฏิบัติ การใช้ข้อมูลและเทคโนโลยี Machine Learning ควรถูกวางบทบาทเป็น “เครื่องมือสนับสนุนการตัดสินใจ” เพื่อคัดกรอง/จัดกลุ่ม วางแผนหลักสูตร ติดตามผล และปรับมาตรการให้เหมาะสมกับบริบทพื้นที่ การดำเนินการดังกล่าวจะช่วยเพิ่มความคุ้มค่าของทรัพยากรภาครัฐ ลดความเหลื่อมล้ำด้านโอกาสในการเข้าถึงการพัฒนา และยกระดับประสิทธิผลของโครงการให้ตอบโจทย์ความแตกต่างของเกษตรกรได้ดียิ่งขึ้น ดังนั้น ข้อค้นพบสำคัญของงานวิจัยนี้คือ ประสิทธิผลของโครงการ Smart Farmer ไม่ได้ขึ้นอยู่กับการออกแบบหลักสูตรเพียงอย่างเดียว แต่ขึ้นอยู่กับความสอดคล้องระหว่างการส่งมอบองค์ความรู้และความพร้อมเชิงโครงสร้างของเกษตรกร การละเลยปัจจัยดังกล่าวนำไปสู่การที่ผลลัพธ์ของโครงการไม่ปรากฏในระดับค่าเฉลี่ย แม้ว่าจะเกิดประโยชน์ในบางกลุ่มก็ตาม
